\documentclass[10pt]{article}

\usepackage[utf8]{inputenc}

\usepackage[T1]{fontenc}

\usepackage{amsmath}

\usepackage{amsfonts}

\usepackage{amssymb}

\usepackage[version=4]{mhchem}

\usepackage{stmaryrd}



\begin{document}

причем дуга $s$ отсчитывается на каждом $L_{k}$ от нек-рой фиксированной точки $z_{k} \in L_{k}$ в положительном направлении; $c(t)=c_{k}=$ const на $L_{k}$. Всегда можно считать $c_{0}=0$, остальные постоянные $c_{k}$ определяются в ходе решения задачи. Если $m=0$, то искомые функции $\varphi$ и $\psi$ голоморфны в $S$. Тогда равенства $\varphi(0)=0, \operatorname{Im} \varphi^{\prime}(0)=0$ обеспечивают единственность решения задачи (1), а необходимые и достаточные условия существования решения



\[

\int_{L}\left(X_{n}+i Y_{n}\right) d s=0, \quad 2 \mu \int_{L}\left(x_{1} Y_{n}-x_{2} X_{n}\right) d s=\operatorname{Re} \int_{L} f d t=0

\]



представляют собой условия статич. равновесия абсолютно жесткого тела.



При $m>0$, как уже было отмечено, $\varphi$ и $\psi$ - многозначные функции специального вида, причем они выражаются через новые искомые голоморфные в области $S$ функции $\varphi^{*}$ и $\psi^{*}$.



\begin{enumerate}

  \setcounter{enumi}{1}

  \item Вторая основная задача: определить упругое равновесие тела по заданным смещениям точек его границы.

\end{enumerate}



Эта задача приводит к граничному условию вида:



\[

x \varphi(t)-t \overline{\varphi^{\prime}(t)}-\overline{\psi(t)}=f(t), \quad t \in L,

\]



где $f=u_{1}+i u_{2}$ - заданная на $L$ функция.



\begin{enumerate}

  \setcounter{enumi}{2}

  \item Основная смешанная задача. Пусть $S-$ конечная односвязная область, ограниченная замкнутым контуром $L ; L=L^{\prime}+L^{\prime \prime}$, где $L^{\prime}$ состоит из конечного числа дуг $L_{1}^{\prime}, \ldots, L_{m}^{\prime}$ контура $L$, к-рые попарно не имеют общих точек; на $L^{\prime}$ заданы внешние напряжения, а на $L^{\prime \prime}$ - смещения. Соответствующие краевые условия можно записать в виде

\end{enumerate}



\[

\gamma(t) \varphi(t)+t \overline{\varphi^{\prime}(t)}+\overline{\psi(t)}=f(t)+c(t), \quad t \in L,

\]



где $f$ - заданная функция точки $t \in L ; \gamma(t)=1$, если $t \in L^{\prime}$, и $\gamma(t)=-x$, если $t \in L^{\prime \prime} ; c(t)=c_{k}=$ const, если $t \in L^{\prime} ; c(t)=0$, если $t \in L^{\prime \prime}$. Постоянные $c_{k}$ (кроме одной, выбираемой как угодно) не задаются заранее и подлежат определению в ходе решения задачи.



\begin{enumerate}

  \setcounter{enumi}{3}

  \item Третья основная задача. На границе области задаются нормальная составляющая вектора смещения и касательная составляющая вектора внешнего напряжения.

\end{enumerate}



Такая задача возникает, напр., в случае соприкасания упругого тела с жестким профилем заданной формы, когда контакт между упругим и жестким телами осуществляется по всей границе. Рассматриваются также другого рода контактные задачи. Все эти задачи также приводятся $\mathbf{K}$ граничным задачам для аналитич. функций.



\begin{enumerate}

  \setcounter{enumi}{4}

  \item Граничные задачи изгиба тонких плас ти нок. K аналогичным граничным условиям приводят задачи изгиба тонких пластинок. Прогиб $w$ серединной поверхности тонкой однородной упругой пластинки, подверженной действию распределенной по ее поверхности нормальной нагрузки интенсивности $q$, удовлетворяет неоднородному бигармонич. уравнению

\end{enumerate}



\[

\Delta \Delta w=q / D,

\]



где $D=E h^{3} / 12\left(1-\sigma^{2}\right)-$ цилиндрич. жесткость; $h-$ толшина пластинки, $E$ - модуль Юнга. Общее решение этого уравнения имеет вид:



\[

w=w_{0}+\tilde{T} q,

\]



где $\tilde{T} q$ - частное решение, к-рое выражается формулой:



\[

\tilde{T} q=(8 \pi D)^{-1} \iint q(\zeta)|\zeta-t|^{2} \ln |\zeta-z| d \zeta_{1} d \zeta_{2},

\]



a $w_{0}$ - общее решение:



\[

w_{0}=\operatorname{Re}(\bar{z} \Phi(z)+\Psi(z))

\]



однородного бигармонич. уравнения



\[

\Delta \Delta w_{0}=0 \text {, }

\]



где $\Phi$ и $\Psi$ - произвольные аналитич. функции в $S$. Если $q$ - многочлен от $x_{1}$ и $x_{2}$, то $\tilde{T} q$ выражается в явной форме. Если $\Phi(0)=0, \Psi(0)=0, \Phi^{\prime}(0)=\overline{\Phi^{\prime}(0)}$, то функции $\Phi$ и $\Psi$ выражаются однозначно посредством заданной бигармонич. функции $w_{0}$.



Решение $w$ уравнения $\Delta \Delta w=q / D$ следует подчинить граничным условиям, соответствующим тому или иному характеру закрепления границы пластинки. В случае заделанной по краям пластинки на границе $L$ области $S$, занятой серединной поверхностью пластинки, должны выполняться условия $w=d w / d n=0$, где $n-$ внешняя нормаль к контуру $L$. Эти два условия можно записать в виде одного комплексного равенства $\partial_{z} w=0$ (на $L$ ). Последнее приводится к виду (1).



Решение этой задачи всегда сушествует, единственно и выражается по формуле:



\[

w\left(x_{1}, x_{2}\right)=D^{-1} \iint S G\left(x_{1}, x_{2}, \zeta_{1}, \zeta_{2}\right) q\left(\zeta_{1}, \zeta_{2}\right) d \zeta_{1} d \zeta_{2},

\]



где $G$ - функция Грина. Для круга $|z|<1$ :



\[

\begin{aligned}

G\left(x_{1}, x_{2}, \zeta_{1}, \zeta_{2}\right) & =2|\zeta-z|^{2} \ln \left(|1-z \bar{\zeta}||\zeta-z|^{-1}\right)- \\

& -\left(1-|z|^{2}\right)\left(1-|\zeta|^{2}\right), \\

z & =x_{1}+i x_{2}, \zeta=\zeta_{1}+i \zeta_{2} .

\end{aligned}

\]



В случае свободной пластинки граничные условия имеют вид:



$\left.\begin{array}{l}\sigma \Delta w+(1-\sigma)\left(w_{x_{1} x_{1}} \cos ^{2} v+w_{x_{2} x_{2}} \sin ^{2} v+w_{x_{1} x_{2}} \sin 2 v\right)=0 \\ \frac{d \Delta w}{d n}+(1-\sigma) \frac{d}{d s}\left(\frac{1}{2}\left(w_{x_{1} x_{2}}-w_{x_{1} x_{1}}\right) \sin 2 v+w_{x_{1} x_{2}} \cos 2 v\right)=0\end{array}\right\}$



где $v-$ угол, составляемый внешней нормалью $n$ с осью $O x_{1}$. Левые части равенств (2) представляют собой соответственно изгибаюший момент и обобщенную перерезываюшую силу, отнесенные к единице длины и действующие на боковой элемент пластинки с нормалью $n$. Граничные условия (2) можно записать в виде:



\[

d\left((3+\sigma) /(1-\sigma) \Phi-z \bar{\Phi}^{\prime}-\bar{\Psi}^{\prime}\right)=g

\]



где $g$ - заданная функция на $L$.



\begin{enumerate}

  \setcounter{enumi}{5}

  \item Плоские стационарные упругие колебан и я. Когда решения системы уравнений динамики упругой среды ищутся в виде

\end{enumerate}



\[

u=v\left(x_{1}, x_{2}\right) e^{\dot{N} t}, u_{3}=0,

\]



где $v$ - частота колебания, для $v$ получаетсяя формула



\[

v=\partial_{\bar{z}}\left(w_{1}+i w_{2}\right)

\]



(предполагается, что внешние силы равны нулю, $X_{\alpha}=0$ ), где $w_{1}$ и $w_{2}$ - произвольные решения уравнений:



\[

\begin{gathered}

\Delta w_{1}+a^{2} v^{2} w_{1}=0, \Delta w_{2}+b^{2} v^{2} w_{2}=0, \\

a^{2}=\rho(\lambda+2 \mu)^{-1}, b^{2}=\rho \mu^{-1} .

\end{gathered}

\]



Поле напряжений выражается по формулам



\[

\begin{gathered}

X_{11}+X_{22}=-(\lambda+\mu) \alpha^{2} v^{2} w_{1}, \\

X_{11}-X_{22}+2 i X_{22}=4 \mu \partial_{\bar{z}}^{2}\left(w_{1}-i w_{2}\right), \\

X_{33}=-2^{-1} a^{2} v^{2} w_{1} .

\end{gathered}

\]



Общее решение уравнения



\[

\Delta w+k^{2} w=0, k^{2}=\mathrm{const}

\]



выражается в виде:



\[

w=A_{0} J_{0}(k|z|)+\int_{0}^{1} \operatorname{Re}[z \Phi(z t)] J_{0}(k|z| \sqrt{1-t}) d t,

\]



где $A_{0}$ - произвольная действительная постоянная, $\Phi-$ произвольная аналитич. функция, $J_{0}$ - функция Бесселя 1-го рода нулевого порядка. При помощи формулы (3) выводятся комплексные представления для полей смещений и напряжений при плоском стационарном колебании упругой среды; они могут быть использованы для исследования граничных задач, а также для построения различных





\end{document}